\section{Lead L4: the two kinds of arc data}

\begin{resultenv}{T4}{Refuted, with an exact cutting-and-gluing replacement}{The 1986 tethers are not the proper chords of R1}
The literal identification in L4 is false. The additional arc datum
described for an interior lieu on cote 156-4/032 has an interior end
and an end on the boundary. A chord of R1 has both ends on the boundary.
These are different types of embedded arc, even in an ordinary disk.
The exact elementary relationship is the following: a proper chord
together with a marked interior point splits into two boundary tethers
with a common interior end and disjoint interiors; conversely, two
such tethers with distinct boundary ends glue uniquely to a proper
chord. For families, the pairwise single-crossing and parallel-class
conditions of R1 must still be imposed separately.
\end{resultenv}
\begin{proof}
We first state exactly what reading this uses. On the original scan
156-4/032, in the paragraph dated 12 June, the additional datum is an
equivalence class of arcs joining an interior lieu to $\partial D$.
We inspected the whole scan and an enlarged crop. Michel Hua's project
was consulted as a secondary reading witness; its corresponding
paragraph leaves the additional datum partly unread, so we do not
take the identification of that datum from the project. On 154/035
the proper-chord requirement is explicit:
$\partial d=d\cap\partial\Delta$. Our assertion compares these typed
arc data. It makes no claim that the fragmentary 1986 discussion is
a complete independent axiom system for pseudo-disks.

Here is a complete counterexample to identifying the two kinds of
arc data, rather than a purported model of all unspecified 1986 axioms.
Let $D=[-1,1]^2$, let the interior lieu be the singleton
$A=\{(0,0)\}$, and take
\[
 \gamma:[0,1]\longrightarrow D,\qquad\gamma(t)=(t,0).
\]
The set $A$ is compact and contractible, and $\gamma$ is an embedded
arc from $A$ to $\partial D$, meeting the boundary only at its final
endpoint. It therefore supplies the boundary-joining arc datum
under discussion. But
\[
 \partial\gamma=\{(0,0),(1,0)\},\qquad
 \gamma\cap\partial D=\{(1,0)\},
\]
so it is not a proper chord. By contrast,
$c(t)=(2t-1,0)$ is a proper chord of the same disk and has two
boundary endpoints. The number of arc endpoints in the ambient
boundary is invariant under every homeomorphism of pairs
$(D,\partial D)\to(D',\partial D')$. Thus no such homeomorphism
identifies these arc data. The example has a finite polygonal
description; \texttt{l4\_endpoint\_certificate.py} checks its endpoint
incidences exactly, without floating-point decisions.

For the positive statement, a homeomorphism $[0,1]\to c$ and a marked
point corresponding to $t_0\in(0,1)$ split $c$ into the two intervals
$[0,t_0]$ and $[t_0,1]$. Each has one interior and one boundary end,
and their intersection is precisely the marked point. Conversely the
union of two embedded intervals with exactly that common endpoint,
otherwise disjoint, is homeomorphic to $[0,1]$. If their other ends
are distinct points of $\partial D$ and their remaining points lie
in the interior, this union is a proper chord. These operations
are inverse on the stated data, up to homeomorphism preserving the
marked point and the unordered pair of chosen tethers; ordering the
pair instead corresponds to orienting the chord. Without the marked point,
pairing, and second tether there is no inverse operation: for example
the single $\gamma$ above can be completed to a chord ending at
$(-1,0)$ or at $(0,1)$, with different boundary endpoint data.
Nothing in this operation enforces a restriction on intersections
of different completed chords. R1's additional conditions are
therefore exactly extra hypotheses, not consequences of the tether
description.

The sentence about two interior lieux on 156-4/045 has a separate
reading problem: the visible affirmative clause about a subdivision
conflicts with the following conclusion, which suggests a missing
negation. We retain both readings. Under the affirmative reading
the two lieux have a common subdivision; under the negative reading
they illustrate an additional compatibility obstruction. Neither
reading changes the endpoint type on /032, the counterexample, or
the cutting-and-gluing proof. Hence no mathematical conclusion here
depends on resolving that doubtful negation. The suggested historical
influence between the two folders is not inferred from the repeated
word for an arc.
\end{proof}
